\documentclass[10pt,a4paper]{amsart}
\usepackage{amsmath,amssymb,amsthm,graphicx}
% Retained minimal preamble for the Stacks Project chapter varieties.tex.
% Verbatim excerpt of the chapter preamble: Xy-pic setup (native preamble lines 14-18), the native theorem
% declarations with the mathematical macros (native preamble lines 154-217), and the project own no-op
% definitions of the reference/slogan/history display environments (native wrapper.tex lines 1-3) that the
% counted statement uses for its slogan. Copied mechanically; nothing added, removed or redefined.
% For commutative diagrams we use Xy-pic
\usepackage[all]{xy}

% We use 2cell for 2-commutative diagrams.
\xyoption{2cell}
% Theorem environments.
%
\theoremstyle{plain}
\newtheorem{theorem}[subsection]{Theorem}
\newtheorem{proposition}[subsection]{Proposition}
\newtheorem{lemma}[subsection]{Lemma}

\theoremstyle{definition}
\newtheorem{definition}[subsection]{Definition}
\newtheorem{example}[subsection]{Example}
\newtheorem{exercise}[subsection]{Exercise}
\newtheorem{situation}[subsection]{Situation}

\theoremstyle{remark}
\newtheorem{remark}[subsection]{Remark}
\newtheorem{remarks}[subsection]{Remarks}

\numberwithin{equation}{subsection}

% Macros
%
\def\lim{\mathop{\mathrm{lim}}\nolimits}
\def\colim{\mathop{\mathrm{colim}}\nolimits}
\def\Spec{\mathop{\mathrm{Spec}}}
\def\Hom{\mathop{\mathrm{Hom}}\nolimits}
\def\Ext{\mathop{\mathrm{Ext}}\nolimits}
\def\SheafHom{\mathop{\mathcal{H}\!\mathit{om}}\nolimits}
\def\SheafExt{\mathop{\mathcal{E}\!\mathit{xt}}\nolimits}
\def\Sch{\mathit{Sch}}
\def\Mor{\mathop{\mathrm{Mor}}\nolimits}
\def\Ob{\mathop{\mathrm{Ob}}\nolimits}
\def\Sh{\mathop{\mathit{Sh}}\nolimits}
\def\NL{\mathop{N\!L}\nolimits}
\def\CH{\mathop{\mathrm{CH}}\nolimits}
\def\proetale{{pro\text{-}\acute{e}tale}}
\def\etale{{\acute{e}tale}}
\def\QCoh{\mathit{QCoh}}
\def\Ker{\mathop{\mathrm{Ker}}}
\def\Im{\mathop{\mathrm{Im}}}
\def\Coker{\mathop{\mathrm{Coker}}}
\def\Coim{\mathop{\mathrm{Coim}}}

% Boxtimes
%
\DeclareMathSymbol{\boxtimes}{\mathbin}{AMSa}{"02}

%
% Macros for moduli stacks/spaces
%
\def\QCohstack{\mathcal{QC}\!\mathit{oh}}
\def\Cohstack{\mathcal{C}\!\mathit{oh}}
\def\Spacesstack{\mathcal{S}\!\mathit{paces}}
\def\Quotfunctor{\mathrm{Quot}}
\def\Hilbfunctor{\mathrm{Hilb}}
\def\Curvesstack{\mathcal{C}\!\mathit{urves}}
\def\Polarizedstack{\mathcal{P}\!\mathit{olarized}}
\def\Complexesstack{\mathcal{C}\!\mathit{omplexes}}
% \Pic is the operator that assigns to X its picard group, usage \Pic(X)
% \Picardstack_{X/B} denotes the Picard stack of X over B
% \Picardfunctor_{X/B} denotes the Picard functor of X over B
\def\Pic{\mathop{\mathrm{Pic}}\nolimits}
\def\Picardstack{\mathcal{P}\!\mathit{ic}}
\def\Picardfunctor{\mathrm{Pic}}
\def\Deformationcategory{\mathcal{D}\!\mathit{ef}}
\newenvironment{reference}{}{}
\newenvironment{slogan}{}{}
\newenvironment{history}{}{}
\usepackage{cleveref}
\crefname{originallemma}{original lemma}{original lemmas}
\Crefname{originallemma}{Original lemma}{Original lemmas}
\makeatletter
\newlabel{lemma-degree-in-terms-of-components}{{0AYW}{source}{}{sourcepointer-1}{}}
\newlabel{lemma-degree-in-terms-of-components@cref}{{[originallemma][0][]0AYW}{[1][0][]source}{}{}{}}
\makeatother
\crefname{originalreference}{original reference}{original references}
\Crefname{originalreference}{Original reference}{Original references}
\makeatletter
\newlabel{divisors-lemma-filter-after-modification}{{0AYP}{source}{}{sourcepointer-2}{}}
\newlabel{divisors-lemma-filter-after-modification@cref}{{[originalreference][0][]0AYP}{[1][0][]source}{}{}{}}
\makeatother
\crefname{originallemma}{original lemma}{original lemmas}
\Crefname{originallemma}{Original lemma}{Original lemmas}
\makeatletter
\newlabel{lemma-degree-birational-pullback}{{0AYU}{source}{}{sourcepointer-3}{}}
\newlabel{lemma-degree-birational-pullback@cref}{{[originallemma][0][]0AYU}{[1][0][]source}{}{}{}}
\makeatother
\crefname{originalreference}{original reference}{original references}
\Crefname{originalreference}{Original reference}{Original references}
\makeatletter
\newlabel{modules-lemma-det-ses}{{0B38}{source}{}{sourcepointer-4}{}}
\newlabel{modules-lemma-det-ses@cref}{{[originalreference][0][]0B38}{[1][0][]source}{}{}{}}
\makeatother
\crefname{originallemma}{original lemma}{original lemmas}
\Crefname{originallemma}{Original lemma}{Original lemmas}
\makeatletter
\newlabel{lemma-degree-tensor-product}{{0AYX}{source}{}{sourcepointer-5}{}}
\newlabel{lemma-degree-tensor-product@cref}{{[originallemma][0][]0AYX}{[1][0][]source}{}{}{}}
\makeatother
\crefname{originallemma}{original lemma}{original lemmas}
\Crefname{originallemma}{Original lemma}{Original lemmas}
\makeatletter
\newlabel{lemma-degree-additive}{{0AYS}{source}{}{sourcepointer-6}{}}
\newlabel{lemma-degree-additive@cref}{{[originallemma][0][]0AYS}{[1][0][]source}{}{}{}}
\makeatother
\title{The degree of a locally free module on a proper curve scheme equals the degree of its determinant}
\author{The Stacks Project Authors}
\date{stacks-project:varieties.tex:local-snapshot-20261011; GNU FDL 1.2 or later}
\begin{document}
\maketitle
\noindent\textit{Editorial note.} All delivered mathematical text is the author's own, copied verbatim from the pinned Stacks Project chapter file varieties.tex; no statement or proof step was rewritten, reordered or removed.\par\medskip
\noindent\textit{Editorial note.} The counted span cites six cross-chapter labels, each recorded in assembly.external\_references with its canonical tag as a navigation pointer.\par\medskip
\noindent\textit{Editorial note.} The proof cites no bibliography item inside the counted span.\par\medskip
\noindent\textit{Allowed standard prerequisites (editorial).}\begin{itemize}
\item Existence of a modification making a locally free module filtrable by invertible quotients. Stacks Project, Divisors, Lemma divisors-lemma-filter-after-modification (tag 0AYP). Produces the filtration on which the determinant computation is performed.
\item Additivity and multiplicativity of the degree, and the determinant of a short exact sequence. Stacks Project, Varieties, Lemmas lemma-degree-additive (tag 0AYS) and lemma-degree-tensor-product (tag 0AYX); Modules on Sites, Lemma modules-lemma-det-ses (tag 0B38). Supply the inductive determinant identity and the final degree computation.
\item Birational invariance of the degree and the componentwise description. Stacks Project, Varieties, Lemma lemma-degree-birational-pullback (tag 0AYU) and Lemma lemma-degree-in-terms-of-components (tag 0AYW). Allow replacing X by the modification and reducing to a single proper curve.
\end{itemize}\medskip
\ifdefined\hypertarget\hypertarget{sourcepointer-1}{}\fi
\noindent\textit{Source reference.} Original lemma 0AYW: sources/stacks-chow/tags.txt, source lines 13804--13804. Cross-chapter standard prerequisite cited by the author inside the counted span; it is used as a pointer and its own proof is not part of this item.\par\medskip
\ifdefined\hypertarget\hypertarget{sourcepointer-2}{}\fi
\noindent\textit{Source reference.} Original reference 0AYP: sources/stacks-chow/tags.txt, source lines 13797--13797. Cross-chapter standard prerequisite cited by the author inside the counted span; it is used as a pointer and its own proof is not part of this item.\par\medskip
\ifdefined\hypertarget\hypertarget{sourcepointer-3}{}\fi
\noindent\textit{Source reference.} Original lemma 0AYU: sources/stacks-chow/tags.txt, source lines 13802--13802. Cross-chapter standard prerequisite cited by the author inside the counted span; it is used as a pointer and its own proof is not part of this item.\par\medskip
\ifdefined\hypertarget\hypertarget{sourcepointer-4}{}\fi
\noindent\textit{Source reference.} Original reference 0B38: sources/stacks-chow/tags.txt, source lines 13956--13956. Cross-chapter standard prerequisite cited by the author inside the counted span; it is used as a pointer and its own proof is not part of this item.\par\medskip
\ifdefined\hypertarget\hypertarget{sourcepointer-5}{}\fi
\noindent\textit{Source reference.} Original lemma 0AYX: sources/stacks-chow/tags.txt, source lines 13805--13805. Cross-chapter standard prerequisite cited by the author inside the counted span; it is used as a pointer and its own proof is not part of this item.\par\medskip
\ifdefined\hypertarget\hypertarget{sourcepointer-6}{}\fi
\noindent\textit{Source reference.} Original lemma 0AYS: sources/stacks-chow/tags.txt, source lines 13800--13800. Cross-chapter standard prerequisite cited by the author inside the counted span; it is used as a pointer and its own proof is not part of this item.\par\medskip
\begin{lemma}
\label{lemma-degree-and-det}
Let $k$ be a field, let $X$ be a proper scheme of dimension $\leq 1$
over $k$, and let $\mathcal{E}$ be a locally free
$\mathcal{O}_X$-module of rank $n$. Then
$$
\deg(\mathcal{E}) = \deg(\wedge^n(\mathcal{E})) = \deg(\det(\mathcal{E}))
$$
\end{lemma}

\begin{proof}
By Lemma \ref{lemma-degree-in-terms-of-components} and elementary
arithmetic, we reduce to the case of a proper curve.
Then there exists a modification $f : X' \to X$ such that
$f^*\mathcal{E}$ has a filtration whose successive
quotients are invertible modules, see
Divisors, Lemma \ref{divisors-lemma-filter-after-modification}.
By Lemma \ref{lemma-degree-birational-pullback} we may work on $X'$.
Thus we may assume we have a filtration
$$
0 = \mathcal{E}_0 \subset \mathcal{E}_1 \subset \mathcal{E}_2 \subset
\ldots \subset \mathcal{E}_n = \mathcal{E}
$$
by locally free $\mathcal{O}_X$-modules
with $\mathcal{L}_i = \mathcal{E}_i/\mathcal{E}_{i - 1}$ is invertible.
By Modules, Lemma \ref{modules-lemma-det-ses} and induction we find
$\det(\mathcal{E}) = \mathcal{L}_1 \otimes \ldots \otimes \mathcal{L}_n$.
Thus the equality follows from Lemma \ref{lemma-degree-tensor-product}
and additivity (Lemma \ref{lemma-degree-additive}).
\end{proof}

\end{document}
